\chapter{Going to Production}
\label{ch:production}\index{production}\index{secure boot}\index{OTA}

Everything to this point has been about making the chip \emph{do} things in Ada. A
product needs more: it has to boot only software you trust, survive a hung task,
take a field update, and not hand its secrets to anyone with a USB cable. This
project's boot path was built for development (it is small, readable and entirely
yours, Chapter~\ref{ch:bootloader}), and it deliberately leaves the production
hardening out. This chapter is honest about that line: what is here, what is not,
and what each missing piece would take. None of the gaps are accidents; the
bootloader's own design notes list secure boot, flash encryption and OTA under an
explicit ``drop'' heading. Knowing precisely where the edge is, is itself part of
shipping safely.

\section{What is on the flash}\index{partition table}

Three things are written, at fixed offsets:

\begin{center}
\begin{tabular}{l l}
\toprule
\textbf{Offset} & \textbf{Contents} \\ \midrule
\texttt{16\#0\#}       & the 2nd-stage bootloader \\
\texttt{16\#8000\#}    & the partition table \\
\texttt{16\#1\_0000\#} & the application image \\
\bottomrule
\end{tabular}
\end{center}

The partition table is a single-app \emph{factory} layout: one application
partition, plus the conventional \texttt{nvs} and \texttt{phy\_init} data regions.
There is one slot for one image. That is the first thing a field-updatable product
changes, and \S\ref{prod:ota} returns to it.

\section{What the bootloader trusts}

When the 2nd-stage loader reads the application header it checks \emph{one} thing:
the \texttt{16\#E9\#} magic byte that marks a valid image. It then reads the segment
count and entry point and runs it. It does not verify the image's checksum (the
\texttt{esp\_elf2image} tool writes one), does not check the appended SHA-256, and
does not (because there is none) check a signature. In development that is
exactly right: it makes flashing instant and the loader trivial to read. In the
field it means the device runs whatever is at offset \texttt{16\#1\_0000\#}, full
stop.

\section{The four gaps, and what each costs to close}

\subsection{Secure boot}\index{secure boot}

\emph{Today:} none. There is no image signing, no verification, and the eFuse
secure-boot bits (which exist in the register definitions) are never read or
burned.

\emph{To close it:} sign the image at build time (RSA-PSS or ECDSA: the pure-Ada
primitives from the TLS work, Chapter~\ref{ch:tls}, already implement both), have
the bootloader verify the signature against a public key before jumping, and burn
the secure-boot eFuse so the mask ROM enforces the same on the 2nd-stage loader
itself. The cryptography is in the repo; the missing parts are the verification step
in \texttt{boot\_main} and a one-time, irreversible eFuse burn.

\subsection{Flash encryption}\index{flash encryption}

\emph{Today:} none. Code and data sit on flash in the clear, readable by anyone
who can dump the chip.

\emph{To close it:} enable the hardware flash-encryption eFuse and key so reads are
transparently decrypted by the MSPI. This one is mostly an eFuse-and-key procedure
rather than Ada code, but it interacts with everything else (an encrypted image must
still verify, an OTA must write ciphertext), so it is best designed in from the
start rather than bolted on.

\subsection{Over-the-air update}\label{prod:ota}\index{OTA}

\emph{Today:} none. The single factory partition has nowhere to stage a new image,
and there is no download-and-flash logic: updating means re-flashing over USB.
The \emph{transport} exists (the W5500 stack and pure-Ada TLS can fetch an image
over HTTPS, Chapters~\ref{ch:networking} and~\ref{ch:tls}); the \emph{mechanism}
does not.

\emph{To close it:} move to an A/B partition layout (two app slots plus an
\texttt{otadata} selector), write the downloaded image into the inactive slot,
verify it (secure boot above), flip the selector, and reboot, with a rollback if
the new image fails to confirm itself healthy. An anti-rollback counter in eFuse
stops a downgrade attack. This is the largest of the four in new code, and it leans
on the other three to be safe; Chapter~\ref{ch:ota} works the design out in full.

\subsection{Locking the debug port}\index{JTAG}

\emph{Today:} the USB-serial-JTAG is enabled, which is how every example in this
book is flashed and debugged. In the field that same port reads RAM and flash.

\emph{To close it:} burn the \texttt{DIS\_PAD\_JTAG} eFuse on a finished unit. It is
one irreversible write, and the project provides no script for it precisely because
it is irreversible: you want to run it deliberately, on a unit you are done
debugging, not as a build step.

\section{Staying alive: watchdogs and resets}\index{watchdog}

There is a second production concern that is not about secrecy but about
\emph{liveness}. The bootloader and startup deliberately \emph{disable} the
watchdogs (the timer-group WDT, the RTC WDT, and the RTC super-watchdog, which is
set to auto-feed) so that a paused debugger or a long elaboration never resets the
chip mid-development. The consequence is that a hung task in the field will hang
\emph{forever}: nothing resets it. A product wants the opposite: a watchdog the
healthy application feeds from its main loop, so that a hang becomes a reset and a
recovery. Re-enabling one of the watchdogs and feeding it is a small, local change,
but it is a change you must make deliberately; the default here is tuned for the
bench, not the field. A brownout detector (unused here) is the matching
analogue for power: worth enabling where supply quality is uncertain.

\section{Where the line sits}

None of this makes the boot path ``insecure''; it makes it \emph{development
grade}, which is the right default for a project whose whole point is to be read,
modified and flashed a hundred times a day. It is well suited to prototyping, to
physically secured deployments, and to learning how the chip actually boots. It is
not, as it stands, suited to a connected device in a hostile environment or to a
fleet that must be updated in the field, and the gap to that is the concrete,
bounded list above, most of it buildable from primitives this project already
contains. The honest summary is the most useful one: the foundation is yours and
you can read every line of it, which is exactly what makes adding the production
layer a matter of engineering rather than of trusting a black box.
